\chapter{Case Study: Storm Surge in Hamburg}
\label{ch:casestudy}

\section{Scenario Overview}
\label{sec:cs:overview}

This thesis uses one concrete case to study evacuation coordination. The case is
a predicted storm surge in Hamburg that forces the evacuation of a nursing home
on the bank of the river Elbe. This section describes the hazard, the site, and
the coordination task. The sections that follow then describe the organizations
involved, the ten steps they go through, and the disruptions used to test them.

\subsection{The Hazard}
\label{sec:cs:overview:hazard}

The Elbe is tidal along its whole course through Hamburg. A storm over the North
Sea can push water up the estuary and raise the water level in the middle of the
city. Hamburg has lived with this risk for a long time, and the risk is expected
to grow with climate change \parencite{vonstorch2008}. Risk studies for the
German Bight and for the Elbe estuary model how such extreme surges develop and
what areas they reach \parencite{oumeraci2015a,wahl2015a}. The case study takes
this hazard as given. The thesis does not model the flood itself. It models what
people and organizations have to do once a surge is forecast.

The important property of the hazard, for this thesis, is the warning time. A
storm surge is not a sudden event. It is forecast hours in advance. This gives
responders a window to act, but the window is fixed and it closes. That makes
the scenario a planning problem under a deadline, which is what the thesis wants
to study.

\subsection{The Site}
\label{sec:cs:overview:site}

The site is the Augustinum, a senior care facility at Elbchaussee 374, 22609
Hamburg. It sits close to the river in the west of the city. The scenario
assumes 100 residents who all have to leave before the water arrives.

The residents are the reason this case is hard. They are old. Many of them
cannot walk far, cannot drive, and cannot be told to leave on their own.
Evacuating them is not a matter of warning people. It is a transport and care
operation. Buses have to be requested, staffed, and routed. Shelters have to be
found that can take a large group at short notice. Medical staff have to travel
with the convoy. Research on evacuation planning for coastal cities shows that
moving elderly populations is the part that dominates the logistics
\parencite{yin2024}, and evacuation models that ignore differences in mobility
underestimate how long the operation takes \parencite{alqahtani2025a}.

Here the thesis also meets a gap in the literature. The German storm-surge work
cited above studies the hazard, the risk, and the damage. It does not study the
evacuation logistics of vulnerable groups on the Elbe coast. The elderly
evacuation work is set in Shanghai and New York, not in Germany. So there is no
ready-made model of this exact operation to reuse, which is one reason the case
is worth building.

\subsection{The Coordination Task}
\label{sec:cs:overview:task}

The task is to move 100 residents from the Augustinum to school buildings that
serve as emergency shelters, before the forecast high water, and to bring them
there safely.

No single organization can do this alone. The city crisis unit knows which
facilities are at risk. The nursing home knows how many residents it has and
what care they need. The schools know how many beds they can offer. The public
transport operator owns the buses. The police can escort a convoy. The medical
service can send a team for the trip. Each of these actors holds information the
others do not have, and each controls resources the others cannot use directly.
The plan only exists once they talk to each other.

Two properties of the case made it a natural place to try an agentic
architecture. First, the actors are heterogeneous and hold local information,
so coordination has to be done by communication and not by a central planner
who already knows everything. Second, the space of things that can go wrong is
large. A road can flood, a shelter can lose capacity, a bus can break down, and
these can happen together and in the middle of the operation. A pre-written
script covers the cases its author thought of. Agents that reason over a shared
database can, in principle, also answer a case nobody wrote down. That promise
is what this thesis set out to test, and Chapter~\ref{ch:eval} describes how.

One boundary belongs here, at the start. The case offers a large disruption
space in principle. The scenarios that were actually run did not use it: none
of them ever \emph{required} a change of plan to save the residents
(Section~\ref{sec:crit:gate0}). The thesis therefore answers whether the agentic
system won on this case. It does not answer whether such a system can ever win.

The scenario is deliberately small enough to run many times and large enough to
be realistic. The roster is fixed at two buses and two receiving shelters, with
one police escort and one medical team. Each bus has 50 seats, so the 100
residents do not fit in a single trip. Each shelter declares 60 beds, so the
120 beds cover the residents with only 20 to spare, and neither shelter on its
own can take everyone.

These numbers are chosen, not incidental. They are what make the allocation a
real decision. A single bus and a single shelter could absorb the whole mission
between them, and then there would be nothing to allocate and nothing to
compare. At two and two, a plan has to split the population across the shelters
and it has to schedule more than one run.

Bus-2 is declared a reserve. The intended plan is that bus-1 makes two
sequential runs and bus-2 stands by for a breakdown. The two systems read this
declaration differently. The baseline sends bus-2 out in parallel from its first
schedule. The agentic transport dispatcher holds it back unless one of four
named reasons releases it. Both sides have the same fleet, so this is a
difference in decisions and not in capability. Section~\ref{sec:val:reserve}
describes the mechanism.

% =====================================================================
\section{Organizational Agents}
\label{sec:cs:agents}

Every agent in this thesis stands for a \emph{person in an organizational role}.
There is a coordinator at a desk, a transport dispatcher, two bus drivers, a
police officer, a paramedic, a head of a nursing home and two head teachers.
Vehicles are not agents. A bus, a patrol car and an ambulance are passive
resources that a person drives. This mirrors how civil protection in Hamburg
actually works: agencies coordinated by people. It also keeps the model honest.
An agent can only know what the person in that role could see from where they
sit.

The case has ten agents in eight organizations
(Table~\ref{tab:cs:roster}).
% TODO(verify): the count of eight organisations is inferred from the roster
% (BIS, ZKD, Hochbahn, Polizei, DRK, Augustinum, and the two schools). No pack
% document lists the eight in one place.
Each agent has a fixed identity with a name, a role, a work address and an
address book of the contacts it may message. The names are placeholders, not
real people.

\begin{table}[htbp]
\centering
\caption{The agent roster. Six agents reason with a language model. Four are
scripted responders.}
\label{tab:cs:roster}
\small
\begin{tabularx}{\textwidth}{l>{\raggedright\arraybackslash}X>{\raggedright\arraybackslash}Xl}
\toprule
Agent & Person and role & Organization & Reasoning \\
\midrule
BIS & Dr.\ Andrea Schneider, Lagedienstleiterin (duty head of the situation centre) & Behörde für Inneres und Sport, Hamburg's interior authority & LLM \\
Hochbahn & Markus Weber, Disponent (dispatcher) in the traffic control centre & Hamburger Hochbahn AG, the public transport operator & LLM \\
Bus-Driver-1 & Hans Müller, drives bus-1 (50 seats, primary) & Hamburger Hochbahn AG & LLM \\
Bus-Driver-2 & Petra Schmidt, drives bus-2 (50 seats, reserve) & Hamburger Hochbahn AG & LLM \\
Polizei & Susanne Krüger, patrol officer, drives patrol car PC-21-1 & Polizei Hamburg, station PK~21 & LLM \\
DRK & Jonas Brandt, paramedic, drives ambulance DRK-T1 & Deutsches Rotes Kreuz (German Red Cross), Altona & LLM \\
ZKD & Dr.\ Thomas Richter, deputy head & Zentraler Katastrophendienststab, the city's central disaster staff & scripted \\
Augustinum & Stefan Becker, Heimleiter (head of the home) & Augustinum, the origin facility & scripted \\
SC1 & Maria Hoffmann, head teacher & Schule Dohrnweg, Altona (60 beds) & scripted \\
SC2 & Klaus Werner, head teacher & Schule Marienthal, Wandsbek (60 beds) & scripted \\
\bottomrule
\end{tabularx}
\end{table}
% TODO(verify): BIS title. PROJECT_BRIEFING and ADR-0008 say
% "Lagedienstleiterin"; agents/bis.md identity says "Leiterin BIS Lagezentrum".

The split between reasoning and scripted agents follows one rule. An agent
that makes no decision in the case gets no language model. The disaster staff
ZKD relays the alarm, books the reports, and countersigns a stand-down. It never
has to choose between options, so it is scripted. The two schools only report
their free beds and confirm who arrived. The nursing home reports its residents
and tells the coordinator when the building is empty. All four answer
mechanically and fast.

The six reasoning agents are the ones that face real choices. BIS is the only
coordinator that decides. It chooses the allocation, asks for transport and
escorts, and is the only agent that may call off the mission. Hochbahn decides
how the fleet is used: which bus runs, in what order, and when the reserve is
released. The police officer and the paramedic are not desk roles. Each of them
drives their own vehicle as part of the convoy, and each plans their own legs.
The two bus drivers do the same for their buses.

The workflow baseline keeps this roster with one change. A single scripted
Incident Commander replaces the two coordinating agents, BIS and Hochbahn. All
other agents are shared by both systems. No coordination decision on the
baseline side therefore comes from a model. Chapter~\ref{ch:workflow} describes
the baseline, and Chapter~\ref{ch:agentic} describes how the agents reason.

% =====================================================================
\section{Coordination Storyline}
\label{sec:cs:storyline}

The mission follows a fixed storyline of ten steps. The storyline says
\emph{what} has to happen and who is responsible. It does not say how each step
is decided. That is left to the coordinator, and it is where the two systems
differ. Figure~\ref{fig:cs:storyline} shows the steps as messages between the
organizations.

\begin{enumerate}
  \item \textbf{Identification.} ZKD receives the flood alert and relays it to
  BIS as an alarm. BIS now knows that a facility in the flood zone has to be
  cleared, and by when.
  \item \textbf{Evacuation assessment.} BIS asks the Augustinum how many
  residents must leave, and passes on the deadline. The home confirms 100.
  \item \textbf{Resource allocation.} BIS asks both schools for free beds. Each
  offers 60. The coordinator now decides how many residents go to each school.
  \item \textbf{Route planning.} No coordinator draws routes. Each driver agent
  routes its own legs over real roads, using OpenRouteService. The coordinators
  work with addresses and times.
  \item \textbf{Logistics.} BIS asks Hochbahn for buses. Hochbahn writes the run
  schedule and sends it back to BIS as a transport plan. By default bus-1 makes
  two sequential runs and bus-2 stays in reserve. If the schedule does not fit
  the deadline, Hochbahn may release the reserve or ask BIS for more time. It
  does not decide the mission alone.
  \item \textbf{Police coordination.} BIS orders an escort. The police officer
  drives the patrol car at the head of the convoy.
  \item \textbf{Medical support.} BIS orders medical cover. The paramedic drives
  the ambulance at the tail of the convoy. Both escorts cover every run and stand
  down only after the last one.
  \item \textbf{Evacuation execution.} The residents board and the bus leaves.
  The driver tells the home that the bus has departed and how many are aboard. A
  bus that only starts boarding does not count. When everyone has left, the home
  reports the building as cleared. If the deadline passes first, it reports the
  building as not cleared, with the number still inside. It sends this report
  once, either way.
  \item \textbf{Arrival and handover.} At the school, the driver hands the
  residents over and the school confirms the intake. The driver also reports the
  drop-off to Hochbahn, which tells BIS that the run is complete.
  \item \textbf{Mission report.} Once every run is complete, BIS sends its report
  to ZKD, which closes the evacuation record.
\end{enumerate}

Steps 8 and 9 repeat for every bus run. The storyline also fixes no split of
residents between the schools. An earlier project brief assumed 50 residents
to each school, one bus each. That is no longer the model. How many go where,
and on which run, is the coordinator's decision.

% The ten-step coordination storyline as message flow between the
% organisations, plus the conditional stand-down. Current design (ADR-0009 to
% ADR-0013), not the pre-ADR-0011 bus-driver sequence diagram.
% Sources: thesis-pack/04-architecture-agentic/agents/bis.md (action table),
% adr/0011 (transport_plan), adr/0012 (facility cleared), adr/0013 (shelter
% edge), adr/0009 (stand-down). Input from inc/casestudy.tex.
\begin{figure}[htbp]
\centering
\begin{tikzpicture}[x=1.0cm, y=1.0cm,
  llm/.style={draw, rounded corners=2pt, fill=blue!8, align=center,
    font=\scriptsize, text width=1.6cm, minimum height=0.9cm, inner sep=2pt},
  scr/.style={draw, rounded corners=2pt, fill=gray!12, align=center,
    font=\scriptsize, text width=1.6cm, minimum height=0.9cm, inner sep=2pt},
  life/.style={densely dashed, gray!60},
  msg/.style={-{Stealth[length=1.6mm]}, thick},
  alt/.style={-{Stealth[length=1.6mm]}, dashed, gray!70!black},
  lbl/.style={above, font=\scriptsize, fill=white, inner sep=1pt, yshift=-1pt},
  note/.style={draw, rounded corners=2pt, fill=orange!8, align=center,
    font=\scriptsize, text width=1.7cm, inner sep=2pt},
]
  % Lanes: x position, style, header text.
  \foreach \x/\s/\t in {
      0/scr/{\textbf{ZKD}\\command},
      2/llm/{\textbf{BIS}\\coordinator},
      4/scr/{\textbf{Augustinum}\\nursing home},
      6/scr/{\textbf{SC1, SC2}\\shelters},
      8/llm/{\textbf{Hochbahn}\\transport},
      10/llm/{\textbf{Bus drivers}\\bus-1, bus-2},
      12/llm/{\textbf{Polizei}\\escort},
      14/llm/{\textbf{DRK}\\ambulance}} {
    \node[\s] at (\x,0) {\t};
    \draw[life] (\x,-0.45) -- (\x,-11.7);
  }

  % 1 Alarm
  \draw[msg] (0,-1.0) -- node[lbl] {1 alarm} (2,-1.0);
  % 2 Assessment
  \draw[msg] (2,-1.6) -- node[lbl] {2 residents?} (4,-1.6);
  \draw[msg] (4,-2.1) -- node[lbl] {100, deadline} (2,-2.1);
  % 3 Resource allocation
  \draw[msg] (2,-2.7) -- node[lbl, pos=0.3] {3 free beds?} (6,-2.7);
  \draw[msg] (6,-3.2) -- node[lbl, pos=0.7] {60 + 60} (2,-3.2);
  % 4 Route planning: done by the drivers, not by a coordinator
  \node[note] at (10,-3.95) {4 route legs on real roads};
  % 5 Logistics
  \draw[msg] (2,-4.7) -- node[lbl, pos=0.35] {5 request buses} (8,-4.7);
  \draw[msg] (8,-5.2) -- node[lbl, pos=0.65] {transport plan} (2,-5.2);
  \draw[msg] (8,-5.7) -- node[lbl] {runs} (10,-5.7);
  % 6 Police escort, 7 medical support
  \draw[msg] (2,-6.3) -- node[lbl, pos=0.85] {6 escort} (12,-6.3);
  \draw[msg] (2,-6.8) -- node[lbl, pos=0.9] {7 medical} (14,-6.8);
  % 8 Execution, repeated per run
  \draw[msg] (10,-7.4) -- node[lbl] {8 departed, $n$ aboard} (4,-7.4);
  \draw[msg] (4,-7.9) -- node[lbl] {facility cleared} (2,-7.9);
  % 9 Arrival and handover, repeated per run
  \draw[msg] (10,-8.5) -- node[lbl] {9 drop-off} (6,-8.5);
  \draw[msg] (10,-9.0) -- node[lbl] {drop-off} (8,-9.0);
  \draw[msg] (8,-9.5) -- node[lbl, pos=0.6] {transport complete} (2,-9.5);
  % 10 Mission report
  \draw[msg] (2,-10.1) -- (0,-10.1);
  \node[font=\scriptsize, anchor=west] at (2.1,-10.1) {10 mission complete};
  % Conditional stand-down (ADR-0009), exercised by no scenario
  \draw[alt] (2,-10.8) -- (0,-10.8);
  \node[font=\scriptsize, anchor=west, text=gray!50!black] at (2.1,-10.8)
    {stand-down only: request abort};
  \draw[alt] (0,-11.3) -- (2,-11.3);
  \node[font=\scriptsize, anchor=west, text=gray!50!black] at (2.1,-11.3)
    {countersigned, then recall of transport and escorts};

  % Legend
  \node[llm, text width=0.6cm, minimum height=0.35cm] at (9.6,-12.15) {};
  \node[font=\scriptsize, anchor=west] at (10.0,-12.15) {LLM-driven};
  \node[scr, text width=0.6cm, minimum height=0.35cm] at (12.0,-12.15) {};
  \node[font=\scriptsize, anchor=west] at (12.4,-12.15) {scripted};
\end{tikzpicture}
\caption[The coordination storyline]{The coordination storyline as message
flow between the organisations of the case. Numbers give the storyline step.
Steps 8 and 9 repeat for every bus run. Step 4 is not a message: each driver
routes its own legs over real roads. The dashed branch is the conditional
stand-down, which no evaluated scenario triggered. In the workflow baseline, one
Incident Commander takes the place of the two coordinator lanes, BIS and
Hochbahn. All other lanes are shared by both systems.}
\label{fig:cs:storyline}
\end{figure}


\paragraph{The stand-down.} One further path exists outside the ten steps. If
BIS finds that no recovery is possible, it asks ZKD to countersign a stand-down.
ZKD authorizes it without conditions, and BIS then recalls the transport, the
police and the medical team. The path keeps the decision with BIS and the
formal authority with ZKD. No evaluated scenario triggered it.

\paragraph{The same storyline in the baseline.} The baseline's Incident
Commander walks the same steps in four phases: \emph{intake} covers steps 1 to
3, \emph{transport allocation} steps 4 and 5, \emph{Befehlsgebung} (issuing
orders) steps 5 to 8, and the \emph{mission report} step 10. Because it
replaces both BIS and Hochbahn, the messages between those two disappear on the
baseline side. Section~\ref{sec:wf:steps} gives the details.

% =====================================================================
\section{Disruption Scenarios}
\label{sec:cs:disruptions}

The storyline describes a mission where nothing goes wrong. The evaluation
disturbs it. Every run begins with the alarm, and every run ends at the forecast
high water. In all scenarios the flood reaches the zone at 12:00. A resident
who reaches a shelter after that time does not count as rescued.
% TODO(verify): 12:00 is sourced from the baseline README and the S07/S08
% rubric wording "by 12:00". Confirm every agentic scenario file declares the same.

Seven scenarios were evaluated. They fall into two groups, which the evaluation
calls the two arms (Table~\ref{tab:eval:scenarios}). The \emph{anticipated} arm
has two scenarios that were known while the systems were built. The baseline has
a written branch for them. The \emph{novel} arm has five scenarios that were
supplied by the domain expert, Rüdiger Höfert, and revealed only after the
baseline had been frozen on 2026-09-10. He wrote them from a brief that named
the \emph{kind} of disruption wanted, compound and unfolding mid-mission, but
not its content. Section~\ref{sec:eval:firewall} describes this firewall. A
third anticipated scenario, a shelter closure, had no agentic counterpart. It
was retired from scoring, and its number, S03, is never reused.

The paragraphs below describe what happens in each scenario and what a sensible
coordinator would do. The checks that turn these expectations into scores are
in Section~\ref{sec:eval:heldout}.

\paragraph{S01, closed loop.} Nothing goes wrong. The two runs go out and land
as planned. This is the undisrupted control. Every other scenario is read
against it.

\paragraph{S02, road closure.} A road on the transport corridor closes while the
evacuation is under way. On the agentic side the closure is reported to the
transport side, and the convoy has to be routed around it. The baseline's
Incident Commander has no routing command. In its version of the scenario the
closure is therefore written as a school becoming unreachable, and the
commander has to re-allocate. The two versions are not the same test. The
baseline version also ends with a lower feasibility ceiling, 60 instead of 100.
Chapter~\ref{ch:validity} discusses what this does to the comparison.

\paragraph{S04, medical emergency in flight.} A resident collapses on a bus that
is already moving. The sensible answer is to keep the bus on its way, send the
ambulance to meet it at the school it is heading for, and tell ZKD. The bus
should be neither diverted nor held. The simulator does not model
the patient, so the case is judged on the decisions, not on the survival count.

\paragraph{S05, unverified persons report.} An emergency call to 112 reports
more people at the home than the plan counts. The report is not confirmed, and
it is later taken back. The sensible answer is to check the report with the
police or ZKD, to keep the running evacuation untouched, and to commit the
reserve bus as a hedge in case the report is true. Once the report is withdrawn, nothing new
should still be on its way.

\paragraph{S06, conflicting evacuation orders.} The district office, the
Bezirksamt, issues an order that contradicts the running plan. The sensible
answer is to answer the home quickly, to escalate the conflict to ZKD, and to
keep the full evacuation going.

\paragraph{S07, reduced bus capacity.} Bus-1 is reported as degraded, and it
may carry only 25 residents instead of 50. The resource is degraded, not lost.
The sensible answer is to load bus-1 with no more than 25, to bring in the
reserve bus, and to keep everyone planned. The simulator does not enforce the
25-seat limit. The scenario file says so in advance, and the case is judged on
the decisions.

\paragraph{S08, partial shelter capacity loss.} A third party reports that
SC1 can now take only 40 residents instead of 60. The sensible answer is to
book no more than 40 into SC1 and to move the rest to SC2. The school itself
still declares 60 beds, because the report is only a message. The case is
therefore judged on where the orders sent people.

\paragraph{What the disruptions have in common.} Every disruption arrives as a
\emph{message}. None of them changes the declared world. The school keeps its
beds, the bus keeps its seats, and the road network stays as it was. A
coordinator that ignores the message can therefore still bring everyone to
safety. This is a property of the case as it was built, and it matters for
reading every result that follows. Section~\ref{sec:crit:gate0} draws the
consequence.

% =====================================================================
\section{Assumptions and Simplifications}
\label{sec:cs:assumptions}

The case is a model, and a model leaves things out. The main simplifications
are listed here. They apply to both systems.

\begin{itemize}
  \item \textbf{Residents are a count.} The 100 residents are a number, not
  agents. They have no behaviour of their own and no subgroups, for example by
  mobility or care level.
  \item \textbf{The flood is a deadline.} The surge itself is not simulated.
  Each zone has one flood arrival time. A resident who is not sheltered by
  then does not count as rescued.
  \item \textbf{Roads have no traffic.} Travel distance and time come from
  OpenRouteService, with a straight-line fallback when the service is not
  available. There is no congestion and no queueing.
  \item \textbf{Bus seats are not enforced.} A bus can in principle carry more
  than its seats. Shelter beds and the flood deadline \emph{are} enforced, on
  both systems. Residents above a shelter's declared beds do not count as safe.
  \item \textbf{Messages always arrive.} Replies come after a fixed delay. There
  is no message loss and no variable delay.
  \item \textbf{One clock speed.} All runs use the same ratio between simulated
  time and real time, 8.4 simulated seconds per real second.
  \item \textbf{One case.} There is one facility, one city, two buses and two
  shelters. Other sizes and layouts were not run.
\end{itemize}

These choices keep the case small enough to run hundreds of times and to score
by a fixed rule. They also bound what the results can say.
Chapter~\ref{ch:critique} and Chapter~\ref{ch:discussion} return to them.
